\explainedfigure{fig_speed_versus_energy_200_episodes}{0.97}{Mean driving speed and estimated energy use of each model}{Each point is one model, placed by its mean speed (horizontal, metres per second) and its estimated traction energy per kilometre (vertical, watt-hours per kilometre) over the 200 clear-weather test episodes. The energy value is the mechanical work at the wheels of a declared virtual vehicle (no regenerative braking), not a battery measurement.}{The two baselines use less energy (85.9 and 90.7 watt-hours per kilometre) and drive faster (8.09 and 8.07 metres per second) than the proposed model (96.2 watt-hours per kilometre at 7.87 metres per second), because they are less careful. The caution-study control is the cheapest (80.4 watt-hours per kilometre).}{Energy comparisons are only meaningful between models with similar safety; a model that crashes early also uses little energy. Differences are valid as comparisons between models, not as absolute battery figures.}{fig:energy}
